\section{Arithmetic calibration of the unchanged diagonal calculus}
\label{sec:pa-calibration}
This section reproduces the mathematical calibration from the original PA195
report \cite{pa195}, with its machine-specific opening and original command
analysis omitted. It is inherited mathematics, included to make the logical
strength of the new table explicit. The new machine uses exactly this
calculus. The changed command compiler is proved separately in
Section~\ref{sec:commands}. All statements here are ordinary mathematical
proofs of PA-provability claims; no new kernel-checked endpoint is asserted.

\subsection{The numerical calculus and its diagonal replacement}

Write
\[
 p(a,b)=\frac{(a+b)(a+b+1)}2+b,
 \quad I(A,B)=4p(A,B),\quad M(a,b)=4p(a,b)+1,
\]
\[
 U(a,B)=4p(B,a)+2,\qquad E(a,b)=4p(a,b)+3.
\]
The notational convention is $\neg A=I(A,1)$, conjunction and existential
quantification being their classical abbreviations. Code $0$ is truth,
and code $1=M(0,0)$ is the designated contradiction. The quantifier
constructor has the transposed Cantor arguments displayed above.

The diagonal calculus $\mathcal C_\Gamma$ retains the six PA196 heads,
the dominated propositional head, modus ponens, and closed
generalization. Its guarded gate is
\begin{align*}
 \Gamma(B)&=I(B,U(B,B)),\\
 V(a,b,c)&=I(M(a,b),U(c,M(a,b))),\\
 P(a,b,c)&=I(c,M(c,b)),\\
 S(a,b,c)&=U\bigl(b,
       (M(a,b)\to P(a,b,c))\wedge(M(c,b)\to c)\bigr),\\
 Q(a,b,c)&=U\bigl(c,S(a,b,c)\to\neg P(a,b,c)\bigr),\\
 F_\Gamma(a,b,c)&=\neg\bigl(\Gamma(c)\to(V(a,b,c)\to Q(a,b,c))\bigr).
\end{align*}
The gate is available when $c\ne a,b$; a failed guard emits $0$.
The logical rules and remaining axiom families are exactly those of
\texttt{PA196/Logic/Basic.lean}, with this one replacement of its
\texttt{bundle}.

Every recursive formula child of a nonzero code is smaller than that
code. Every variable label, including a bound label, occurring in a
formula of code $B$ is strictly below $B$. These assertions follow from
the inequalities $a,b\le p(a,b)$ and the four tags, with code zero handled
separately. Consequently $x_B$ is absent from the decoded formula $B$.
The diagonal clause is a valid vacuity instance, although its binder is
chosen from the numerical code of the entire formula.

Classical implicational reasoning extracts from the gate both
$\Gamma(c)$ and the old gate $G(a,b,c)=\neg(V(a,b,c)\to Q(a,b,c))$.
Every diagonal instance $\Gamma(B)$ is available: use $c=B$ and
$a=b=B+1$. Thus all membership-vacuity instances and pointed separation
used in the original proof remain available.

\begin{lemma}[Recovery of the required equality interfaces]
\label{lem:pa-diagonal-interfaces}
The diagonal calculus derives, uniformly in the numerical labels,
\begin{align*}
 E(a,b)&\to\bigl(M(d,a)\leftrightarrow M(d,b)\bigr),\\
 E(a,b)&\to\bigl(M(a,d)\leftrightarrow M(b,d)\bigr),\\
 E(a,b)&\to R(a,b,z)\quad (z\ne a,b),
\end{align*}
where
\[
 R(a,b,z)=U\bigl(z,M(a,z)\leftrightarrow M(b,z)\bigr)
 \wedge U\bigl(z,M(z,a)\leftrightarrow M(z,b)\bigr),
\]
and it derives $\exists x_a R(a,b,z)$ for $z\ne a,b$.
All these derivations are supplied by primitive-recursive proof templates,
with correctness provable in PA.
\end{lemma}

\begin{proof}
The two forward membership rewrites are retained heads. To obtain the
reverse directions, first assume $a\ne b$. Take
$\chi(t)=M(d,t)$ or $M(t,d)$, initially with $d\ne a$, or take the
diagonal family $\chi(t)=M(t,t)$. Put
\[
 C=E(a,b)\to(\chi(b)\to\chi(a)),\qquad
 \varphi=\neg C,\qquad c=\ulcorner\varphi\urcorner.
\]
Here $c$ is the actual numerical code of the displayed formula, so it is
larger than every label in $C$. In particular, all required guards hold.
Apply the old gate at $(a,a,c)$ and abbreviate its $S,P,Q$ by those
letters. We need only its consequence $\neg Q$, not an unproved general
rule for changing bound-variable names.

The second conjunct in $S$ gives
$S\to U(a,M(c,a)\to\varphi)$. Propositionally,
$\varphi$ implies $E(a,b)$, $\chi(b)$, and $\neg\chi(a)$; hence universal
monotonicity gives
\[
 S\to U(a,M(c,a)\to\neg\chi(a)).
\]
The equality atom has disappeared from this quantified consequent.
The unchanged negative-property specialization lemma now gives
\[
 S\to\bigl(M(c,b)\to\neg\chi(b)\bigr).                 \tag{1}
\]
For precision, this specialization uses only the retained witness
$\exists x_a E(b,a)$, direct rewriting under $E(b,a)$, and vacuity of
$M(c,b)\to\neg\chi(b)$ at $a$. The latter follows from
membership vacuity because $a\ne b,c$ and, in the two nondiagonal
families, $a\ne d$. It is exactly the construction
\texttt{negativeSpecialization} in \texttt{PA196/Logic/Reverse.lean};
none of its premises is primitive equality vacuity.

Under $\varphi$, the point condition $P$ gives $M(c,a)$.
Direct element-slot rewriting under $E(a,b)$ then gives $M(c,b)$.
Together with (1), $S$ contradicts the consequence $\chi(b)$ of
$\varphi$. Classical deduction gives
$\varphi\to(S\to\neg P)$. Generalize this theorem at $c$, distribute
the quantifier through implication, and use
$\Gamma(\varphi)=\varphi\to U(c,\varphi)$. The result is
$\varphi\to Q$. Since the gate supplies $\neg Q$, double-negation
elimination gives $C$, as required.

If $a=b$, the reverse rewrite is propositional. When the fixed label
$d=a$, direct rewriting first takes $M(a,b)$ or $M(b,a)$ to $M(b,b)$;
the already proved diagonal case takes $M(b,b)$ to $M(a,a)$.
This covers every label coincidence without a freshness assumption on
the theorem being derived.

For the neighborhood formula choose $z_0=\ulcorner E(a,b)\urcorner$.
The available diagonal clause $E(a,b)\to U(z_0,E(a,b))$ and the two
bidirectional rewrites yield $E(a,b)\to R(a,b,z_0)$.
To change the binder to an arbitrary $z\ne a,b$, only the case
$z\ne z_0$ remains. Let $P(t)$ be either membership-only matrix in
$R$. Bidirectional rewriting gives
$E(z,z_0)\to(P(z_0)\to P(z))$.
The witness $\exists x_{z_0}E(z,z_0)$, universal distribution, and
vacuity of $P(z)$ at $z_0$ give
\[
 U(z_0,P(z_0))\to P(z).
\]
Its antecedent is vacuous at $z$: its nonconstant membership atoms
have indices only $a,b,z_0$. The designated false code $1=M(0,0)$,
which occurs in the Boolean abbreviations, is vacuous at every label
by explosion. Vacuity extends through the Boolean connectives and the
outer quantifier by the retained commutation head.
Generalization and distribution therefore give
$U(z_0,P(z_0))\to U(z,P(z))$.
This proves the claimed neighborhood implication for every allowed $z$.
Finally the unchanged witness head and existential monotonicity give
the neighborhood witness, reversing its two arguments by the
propositional symmetry of $R$ if necessary.

Every operation in this proof is a fixed syntactic proof template,
composition, numerical code formation, or a displayed decidable
label-equality case split. The choices $c=\ulcorner\neg C\urcorner$ and
$z_0=\ulcorner E(a,b)\urcorner$ use no unbounded search. PA verifies the
templates by the same proof-checker induction as their original
counterparts. This establishes the asserted uniformity.
\end{proof}

The PA196 refutation compiler uses primitive equality through these
membership-rewriting, neighborhood, and witness interfaces. In its
logical and arithmetic certificate copies, the uses of equality vacuity
are in the reverse-rewriting and neighborhood modules just replaced.
The old gate, all other heads, and the command stack interface remain
available. Substituting Lemma~\ref{lem:pa-diagonal-interfaces}'s templates
therefore gives a PA-verifiable refutation compiler for
$\mathcal C_\Gamma$. In particular,
\begin{equation}\label{eq:pa-forward-calculus}
 \mathrm{PA}\vdash
 \exists p\,\operatorname{Prf}_{\mathrm{PA}}(p,\ulcorner0=1\urcorner)
 \ \longrightarrow\
 \exists d\,\operatorname{Prf}_{\Gamma}(d,1).
\end{equation}
This is reuse of the published proof-producing construction with explicit
replacement templates, rather than an appeal to standard-model
equiconsistency. The 191-item AX word still pushes nine items in the
same positions as the original nine-item command, so its stack
compiler indices need no alteration.

\subsection{A uniform translation of the numerical calculus into PA}

We now prove the converse proof-theoretic direction. Work first in the
fixed conservative definitional extension of PA containing
$\operatorname{pow2}(n)=2^n$ and its usual arithmetical definition.
Only fixed primitive-recursive functions are added. A fixed
primitive-recursive elimination of these definitions converts the
resulting proofs to ordinary PA proofs.

Use the bit predicate
\[
 B(s,z)\;\equiv\;
 \exists u,v\,[s=(2u+1)2^z+v\ \wedge\ v<2^z].
\]
PA proves the usual digit lemmas directly by induction and division
with remainder. In particular $B(s,z)\to z<s$, since
$s\ge 2^z\ge z+1$, and hence PA proves $\neg B(x,x)$.

Let $T(c)$ be the arithmetic formula obtained by decoding the source
code $c$:
\begin{align*}
 T(0)&=(0=0),& T(M(a,b))&=B(x_a,x_b),\\
 T(E(a,b))&=(x_a=x_b),& T(U(a,A))&=\forall x_a\,T(A).
\end{align*}
For a nonzero implication code the translation is the corresponding
implication. There is one numerical alias:
$I(0,0)=0$. Thus the uniform constructor theorem is a proof of
\[
 T(I(A,B))\leftrightarrow(T(A)\to T(B)),
\]
not an unconditional assertion of syntactic identity. At $A=B=0$ it
uses the fixed proof $(0=0)\leftrightarrow((0=0)\to(0=0))$.
Every other case follows from the definition of $T$.

Arithmetic auxiliary bound variables are represented with de Bruijn
indices, or equivalently renamed above all source labels before
substitution. The source label $c$ is therefore absent among the free
variables of $T(c)$ even though the expanded bit formula has its own
bound variables. Quantifier formation closes the indicated free source
label, so nested occurrences of the same label have their usual
shadowing semantics. No assumption that all numerical labels are
distinct is made.

\begin{lemma}[Uniform finite bit comprehension]
\label{lem:pa-bit-comprehension}
Given the syntax code of any arithmetic formula $\theta(z,\vec u)$,
there is a primitive-recursive procedure producing a PA proof of
\[
 \forall n\,\exists s<2^n\,\forall z<n\,
       \bigl(B(s,z)\leftrightarrow\theta(z,\vec u)\bigr).
                                                        \tag{2}
\]
PA proves that the procedure produces a valid proof for every
well-formed formula input, with the indicated fresh variables.
\end{lemma}

\begin{proof}
Take induction on $n$ with the existential formula in (2) as induction
formula. At zero choose $s=0$. Given a witness $s<2^n$, split by
excluded middle on $\theta(n,\vec u)$. If it holds use $s+2^n$;
otherwise use $s$. The fixed bit lemmas say that the first $n$ bits
are unchanged, that the new $n$th bit is respectively one or zero,
and that the result is below $2^{n+1}$. This is a fixed finite
first-order proof pattern with one PA induction instance.

On syntax codes the procedure performs capture-avoiding substitution,
forms the displayed induction instance, and concatenates this fixed
proof pattern with the fixed arithmetic lemma proofs. It never
evaluates $\theta$. Its syntax operations and a fresh-bound-variable
choice are primitive recursive. The PA proof of correctness is
induction on the syntax operations, followed by direct verification
that the generated induction line is an instance of the PA induction
schema and that each subsequent line follows by the stated logical
rule. Thus this is a statement about proof checking, with an arbitrary
formula code as input, not a uniform satisfaction or reflection
principle for PA. Full PA induction is used: $\theta$ is not restricted
to a bounded or quantifier-free formula.
\end{proof}

We spell out the gate supplier because label collisions matter here.
Fix $a,b,c$ with $c\ne a,b$, let $\phi=T(c)$, and regard $x_a,x_b$
as the free values before the inner quantifier is entered.
If $a\ne b$, apply (2) with bound $x_a+x_b+1$ and matrix
\[
 \bigl(B(x_a,z)\vee z=x_b\bigr)\wedge\phi[z/x_b].
                                                        \tag{3}
\]
The witness $s$ satisfies, for every $z$,
\[
 B(x_a,z)\to(\phi[z/x_b]\to B(s,z)),\qquad
 B(s,z)\to\phi[z/x_b],
\]
and $\phi\to B(s,x_b)$. To pass from the bounded comprehension
formula to the displayed unbounded assertions, use
$B(x_a,z)\to z<x_a$ and the fact that a number below $2^n$ has no
set bit at an index at least $n$.

If the \emph{labels} $a=b$ coincide, the first antecedent inside
$T(S)$ becomes $B(z,z)$, not $B(x_a,z)$ with $x_a$ held fixed.
It is therefore false. Apply (2) instead with bound $x_b+1$ and
matrix
\[
 (z=x_b)\wedge\phi[z/x_b].                              \tag{4}
\]
This witness has the required second implication and distinguished
point condition, and the first implication is vacuous. Equality
substitution in (4) supplies $\phi$ at the distinguished point.
The substitutions in (3)--(4) are capture avoiding; in particular
the inner binder corresponding to $b$ is renamed when the witness
expression depends on the original free value of $x_b$.

In both cases the result is a PA proof of
$\exists x_c\,(T(S)\wedge T(P))$. Freshness of $c$ in $\phi$
ensures that choosing the set witness does not alter the property.
Consequently PA proves $\neg T(Q)$. The guarded membership-vacuity
clause has a first-order proof, and $T(\Gamma(c))$ has a first-order
proof by the code-freshness lemma. Classical propositional reasoning
therefore proves $T(F_\Gamma(a,b,c))$.
When the guard fails the command emits $0$, with its fixed truth proof.

The source false code is translated to $T(1)=B(x_0,x_0)$.
Its negation is a fixed PA theorem, uniformly under all binders.
Thus source negation, conjunction, and existential abbreviations have
PA-provable equivalences to the ordinary classical connectives, even
if a surrounding source quantifier binds label zero. This deals with
the false-code convention independently of the truth-code alias.

The other head suppliers are fixed first-order proof templates:
the witness-elimination head, quantifier commutation, the two equality
congruence heads, an equality witness, and quantifier distribution.
Coincident labels are handled by capture-avoiding instantiation, with
a numerical equality case split if needed. The dominated
propositional head is a classical tautology; its arithmetic validity
does not require the dominance restriction. Generalization and modus
ponens use their ordinary PA proof constructors, together with the
implication-constructor equivalence above.

\begin{lemma}[Syntactic refutation translation]
\label{lem:pa-refutation-translation}
There is a primitive-recursive function $J$ such that
\[
 \mathrm{PA}\vdash\forall d\,
 \left(\operatorname{Prf}_{\Gamma}(d,1)
 \to\operatorname{Prf}_{\mathrm{PA}}
             (J(d),\ulcorner0=1\urcorner)\right).
                                                        \tag{5}
\]
\end{lemma}

\begin{proof}
Decode the source derivation by recursion on its finite proof code.
At an axiom node insert the appropriate supplier just described.
At modus ponens concatenate the two translated proofs and the
compositional implication proof; at generalization append the
corresponding quantifier rule. At a code-zero alias emit or use the
fixed truth equivalence. Recursion on a malformed proof is made total
by an arbitrary fixed output; the correctness statement has the
displayed proof-validity hypothesis.

PA proves the syntactic invariant
$\operatorname{Prf}_{\Gamma}(d,A)\to
  \operatorname{Prf}_{\mathrm{PA}}(J_0(d),T(A))$
by induction on this coded recursion. Constructor children decrease
strictly; proof concatenation, label closing and fresh-variable
renaming have primitive-recursive definitions and their usual
syntactic correctness inductions. The gate case is precisely
Lemma~\ref{lem:pa-bit-comprehension} and the two label cases above.
Append the fixed theorem $\neg B(x_0,x_0)$ at a proof of code one,
then eliminate the fixed definitional extension. This gives (5).

This argument proves one quantified assertion about proof codes
inside PA. It does not merely produce a PA proof separately for each
standard formula and then assert a uniform conclusion. The inductive
verification of the compiler, including the schema supplier, is the
step that makes (5) applicable to nonstandard codes in every model
of PA.
\end{proof}

